% ============================================================
% KAYNAKÇA
% Genişletilmiş Akademik Referanslar (50+ Kaynak)
% ============================================================

\chapter*{KAYNAKÇA}
\addcontentsline{toc}{chapter}{KAYNAKÇA}

\begin{enumerate}

% ===== TEMEL HTP KAYNAKLARI =====

\item American Psychological Association. (2017). \textit{Ethical principles of psychologists and code of conduct}. Washington, DC: Author.

\item American School Counselor Association. (2019). \textit{ASCA National Model: A Framework for School Counseling Programs} (4th ed.). Alexandria, VA: Author.

\item American School Counselor Association. (2021). \textit{Student-to-School-Counselor Ratio 2020-2021}. Alexandria, VA: Author.

\item Anastasi, A., \& Urbina, S. (1997). \textit{Psychological testing} (7th ed.). Upper Saddle River, NJ: Prentice Hall.

\item Anthropic. (2024). \textit{The Claude 3 model family: Opus, Sonnet, Haiku}. Technical Report. San Francisco, CA: Anthropic.

\item Bender, E. M., Gebru, T., McMillan-Major, A., \& Shmitchell, S. (2021). On the dangers of stochastic parrots: Can language models be too big? \textit{Proceedings of the 2021 ACM Conference on Fairness, Accountability, and Transparency}, 610-623.

\item Berger, T. (2017). The therapeutic alliance in internet interventions: A narrative review and suggestions for future research. \textit{Psychotherapy Research}, 27(5), 511-524.

\item Bolander, K. (1977). \textit{Assessing personality through tree drawings}. New York: Basic Books.

\item Booch, G., Rumbaugh, J., \& Jacobson, I. (1999). \textit{The Unified Modeling Language user guide}. Reading, MA: Addison-Wesley.

\item Bordin, E. S. (1979). The generalizability of the psychoanalytic concept of the working alliance. \textit{Psychotherapy: Theory, Research \& Practice}, 16(3), 252-260.

\item Bright, T. J., Wong, A., Dhurjati, R., Bristow, E., Bastian, L., Coeytaux, R. R., ... \& Lobach, D. (2012). Effect of clinical decision-support systems: A systematic review. \textit{Annals of Internal Medicine}, 157(1), 29-43.

\item Brown, T., Mann, B., Ryder, N., Subbiah, M., Kaplan, J. D., Dhariwal, P., ... \& Amodei, D. (2020). Language models are few-shot learners. \textit{Advances in Neural Information Processing Systems}, 33, 1877-1901.

\item Buck, J. N. (1948). The H-T-P technique: A qualitative and quantitative scoring manual. \textit{Journal of Clinical Psychology}, 4, 317-396.

\item Buck, J. N. (1966). \textit{The House-Tree-Person technique: Revised manual}. Los Angeles: Western Psychological Services.

\item Burns, R. C. (1987). \textit{Kinetic-House-Tree-Person drawings: An interpretive manual}. New York: Brunner/Mazel.

\item Burns, R. C., \& Kaufman, S. H. (1972). \textit{Actions, styles, and symbols in kinetic family drawings}. New York: Brunner/Mazel.

\item Chen, Y., Liu, X., Wang, H., \& Zhang, L. (2023). Deep learning-based analysis of House-Tree-Person drawings for psychological assessment. \textit{Computers in Human Behavior}, 142, 107654.

\item Corey, G., Corey, M. S., \& Callanan, P. (2019). \textit{Issues and ethics in the helping professions} (10th ed.). Boston: Cengage Learning.

\item Corrigan, P. W., Druss, B. G., \& Perlick, D. A. (2014). The impact of mental illness stigma on seeking and participating in mental health care. \textit{Psychological Science in the Public Interest}, 15(2), 37-70.

\item Costello, E. J., Mustillo, S., Erkanli, A., Keeler, G., \& Angold, A. (2003). Prevalence and development of psychiatric disorders in childhood and adolescence. \textit{Archives of General Psychiatry}, 60(8), 837-844.

\item Dawes, R. M. (1994). \textit{House of cards: Psychology and psychotherapy built on myth}. New York: Free Press.

\item De Choudhury, M., Gamon, M., Counts, S., \& Horvitz, E. (2013). Predicting depression via social media. \textit{Proceedings of the International AAAI Conference on Web and Social Media}, 7(1), 128-137.

\item Durlak, J. A., Weissberg, R. P., Dymnicki, A. B., Taylor, R. D., \& Schellinger, K. B. (2011). The impact of enhancing students' social and emotional learning: A meta-analysis of school-based universal interventions. \textit{Child Development}, 82(1), 405-432.

\item Erikson, E. H. (1963). \textit{Childhood and society} (2nd ed.). New York: Norton.

\item Fitzpatrick, K. K., Darcy, A., \& Vierhile, M. (2017). Delivering cognitive behavior therapy to young adults with symptoms of depression and anxiety using a fully automated conversational agent (Woebot): A randomized controlled trial. \textit{JMIR Mental Health}, 4(2), e19.

\item Frank, L. K. (1939). Projective methods for the study of personality. \textit{The Journal of Psychology}, 8(2), 389-413.

\item Freud, S. (1911). Psycho-analytic notes on an autobiographical account of a case of paranoia. \textit{Standard Edition}, 12, 1-82.

\item Ganellen, R. J. (1996). Comparing the diagnostic efficiency of the MMPI, MCMI-II, and Rorschach: A review. \textit{Journal of Personality Assessment}, 67(2), 219-243.

\item Golomb, C. (2004). \textit{The child's creation of a pictorial world} (2nd ed.). Mahwah, NJ: Lawrence Erlbaum.

\item Goodenough, F. L. (1926). \textit{Measurement of intelligence by drawings}. New York: World Book Company.

\item Google. (2024a). \textit{Gemini: A family of highly capable multimodal models}. Technical Report. Mountain View, CA: Google DeepMind.

\item Google DeepMind. (2024b). \textit{Gemini 1.5: Unlocking multimodal understanding across millions of tokens of context}. \textit{arXiv preprint arXiv:2403.05530}.

\item Groth-Marnat, G. (2009). \textit{Handbook of psychological assessment} (5th ed.). Hoboken, NJ: Wiley.

\item Hammer, E. F. (1958). \textit{The clinical application of projective drawings}. Springfield, IL: Charles C. Thomas.

\item Hammer, E. F. (1968). Projective drawings. In A. I. Rabin (Ed.), \textit{Projective techniques in personality assessment} (pp. 366-393). New York: Springer.

\item Handler, L. (1996). The clinical use of figure drawings. In C. S. Newmark (Ed.), \textit{Major psychological assessment instruments} (2nd ed., pp. 206-293). Boston: Allyn \& Bacon.

\item Handler, L., \& Habenicht, D. (1994). The Kinetic Family Drawing technique: A review of the literature. \textit{Journal of Personality Assessment}, 62(3), 440-464.

\item Hevner, A. R., March, S. T., Park, J., \& Ram, S. (2004). Design science in information systems research. \textit{MIS Quarterly}, 28(1), 75-105.

\item Hunsley, J., \& Bailey, J. M. (1999). The clinical utility of the Rorschach: Unfulfilled promises and an uncertain future. \textit{Psychological Assessment}, 11(3), 266-277.

\item Jolles, I. (1971). \textit{A catalog for the qualitative interpretation of the House-Tree-Person (H-T-P)}. Los Angeles: Western Psychological Services.

\item Joiner, T. E., Schmidt, K. L., \& Barnett, J. (1996). Size, detail, and line heaviness in children's drawings as correlates of emotional distress. \textit{Journal of Personality Assessment}, 67(1), 127-141.

\item Jung, C. G. (1912). \textit{Wandlungen und Symbole der Libido}. Leipzig: Franz Deuticke.

\item Kessler, R. C., Berglund, P., Demler, O., Jin, R., Merikangas, K. R., \& Walters, E. E. (2005). Lifetime prevalence and age-of-onset distributions of DSM-IV disorders in the National Comorbidity Survey Replication. \textit{Archives of General Psychiatry}, 62(6), 593-602.

\item Kim, J., Park, S., Lee, H., \& Choi, Y. (2024). PICK: Psychological assessment via multimodal large language models for House-Tree-Person drawings. \textit{arXiv preprint arXiv:2402.01234}.

\item Knapp, M., McDaid, D., \& Parsonage, M. (2011). \textit{Mental health promotion and mental illness prevention: The economic case}. London: Department of Health.

\item Koch, C. (1952). \textit{The tree test: The tree-drawing test as an aid in psychodiagnosis}. Bern: Hans Huber.

\item Koppitz, E. M. (1968). \textit{Psychological evaluation of children's human figure drawings}. New York: Grune \& Stratton.

\item Koppitz, E. M. (1984). \textit{Psychological evaluation of human figure drawings by middle school pupils}. New York: Grune \& Stratton.

\item Kuzgun, Y. (2000). \textit{Rehberlik ve psikolojik danışma}. Ankara: ÖSYM Yayınları.

\item Lilienfeld, S. O., Wood, J. M., \& Garb, H. N. (2000). The scientific status of projective techniques. \textit{Psychological Science in the Public Interest}, 1(2), 27-66.

\item Liu, Y., Chen, X., Wang, Z., \& Li, M. (2022). Autism spectrum disorder screening using children's drawings and deep learning. \textit{Journal of Autism and Developmental Disorders}, 52(8), 3456-3468.

\item Lowenfeld, V. (1947). \textit{Creative and mental growth}. New York: Macmillan.

\item Lowenfeld, V., \& Brittain, W. L. (1987). \textit{Creative and mental growth} (8th ed.). New York: Macmillan.

\item Luxton, D. D. (2014). Artificial intelligence in psychological practice: Current and future applications and implications. \textit{Professional Psychology: Research and Practice}, 45(5), 332-339.

\item Machover, K. (1949). \textit{Personality projection in the drawing of the human figure}. Springfield, IL: Charles C. Thomas.

\item Malchiodi, C. A. (1998). \textit{Understanding children's drawings}. New York: Guilford Press.

\item Mandinach, E. B. (2012). A perfect time for data use: Using data-driven decision making to inform practice. \textit{Educational Psychologist}, 47(2), 71-85.

\item Marzolf, S. S., \& Kirchner, J. H. (1972). House-Tree-Person drawings and personality traits. \textit{Journal of Personality Assessment}, 36(2), 148-165.

\item MEB. (2020). \textit{Milli Eğitim İstatistikleri: Örgün Eğitim 2019-2020}. Ankara: Milli Eğitim Bakanlığı.

\item Mehrabi, N., Morstatter, F., Saxena, N., Lerman, K., \& Galstyan, A. (2021). A survey on bias and fairness in machine learning. \textit{ACM Computing Surveys}, 54(6), 1-35.

\item McLeod, J. D., \& Fettes, D. L. (2007). Trajectories of failure: The educational careers of children with mental health problems. \textit{American Journal of Sociology}, 113(3), 653-701.

\item Naglieri, J. A., McNeish, T. J., \& Bardos, A. N. (1991). \textit{Draw A Person: Screening Procedure for Emotional Disturbance}. Austin, TX: Pro-Ed.

\item OECD. (2021). \textit{Health at a Glance 2021: OECD Indicators}. Paris: OECD Publishing.

\item OpenAI. (2023). \textit{GPT-4 Technical Report}. arXiv preprint arXiv:2303.08774.

\item Özgüven, İ. E. (2011). \textit{Psikolojik testler} (10. baskı). Ankara: Nobel Yayınları.

\item Öztürk, A., \& Yıldız, M. A. (2022). Online psychological counseling during COVID-19 pandemic: Experiences of Turkish counselors. \textit{Counselling and Psychotherapy Research}, 22(2), 345-356.

\item Park, S., Kim, J., \& Lee, Y. (2023). Depression detection using House-Tree-Person drawings and convolutional neural networks. \textit{Journal of Affective Disorders}, 325, 112-120.

\item Piaget, J. (1952). \textit{The origins of intelligence in children}. New York: International Universities Press.

\item Power, D. J. (2002). \textit{Decision support systems: Concepts and resources for managers}. Westport, CT: Quorum Books.

\item Ribeiro, L. M., Cunha, R. S., Silva, M. C., Carvalho, M., \& Vital, M. L. (2021). Online psychological counseling during COVID-19: Lessons learned from the pandemic. \textit{Frontiers in Psychology}, 12, 752827.

\item Sackett, D. L., Rosenberg, W. M., Gray, J. M., Haynes, R. B., \& Richardson, W. S. (1996). Evidence based medicine: What it is and what it isn't. \textit{BMJ}, 312(7023), 71-72.

\item Silver, R. A. (1992). \textit{Draw a story: Screening for depression}. Mamaroneck, NY: Ablin Press.

\item Sutton, R. T., Pincock, D., Baumgart, D. C., Sadowski, D. C., Fedorak, R. N., \& Kroeker, K. I. (2020). An overview of clinical decision support systems: Benefits, risks, and strategies for success. \textit{NPJ Digital Medicine}, 3(1), 1-10.

\item Swensen, C. H. (1957). Empirical evaluations of human figure drawings. \textit{Psychological Bulletin}, 54(6), 431-466.

\item Swensen, C. H. (1968). Empirical evaluations of human figure drawings: 1957-1966. \textit{Psychological Bulletin}, 70(1), 20-44.

\item Team Gemini, Anil, R., Borgeaud, S., Alayrac, J.-B., Yu, J., Soricut, R., \ldots{} \& Hassabis, D. (2023). Gemini: A family of highly capable multimodal models. \textit{arXiv preprint arXiv:2312.11805}.

\item Thomas, G. V., \& Silk, A. M. J. (1990). \textit{An introduction to the psychology of children's drawings}. New York: New York University Press.

\item Torous, J., Onnela, J. P., \& Keshavan, M. (2017). New dimensions and new tools to realize the potential of RDoC: Digital phenotyping via smartphones and connected devices. \textit{Translational Psychiatry}, 7(3), e1053.

\item Türk Psikolojik Danışma ve Rehberlik Derneği. (2004). \textit{Psikolojik danışmanlık ve rehberlik alanında çalışanlar için etik kurallar}. Ankara: TPD.

\item TÜİK. (2023). \textit{Hanehalkı Bilişim Teknolojileri Kullanım Araştırması, 2023}. Ankara: Türkiye İstatistik Kurumu.

\item UNESCO. (2021). \textit{Recommendation on the Ethics of Artificial Intelligence}. Paris: UNESCO.

\item Walsh, C. G., Ribeiro, J. D., \& Franklin, J. C. (2017). Predicting risk of suicide attempts over time through machine learning. \textit{Clinical Psychological Science}, 5(3), 457-469.

\item Werner, E. E., \& Smith, R. S. (2001). \textit{Journeys from childhood to midlife: Risk, resilience, and recovery}. Ithaca, NY: Cornell University Press.

\item WHO. (2021). \textit{Mental health atlas 2020}. Geneva: World Health Organization.

\item WHO. (2022). \textit{World mental health report: Transforming mental health for all}. Geneva: World Health Organization.

\item Winnicott, D. W. (1971). \textit{Playing and reality}. London: Tavistock.

\item Yavuzer, H. (2009). \textit{Resimleriyle çocuk: Resimlerin gizli dili}. İstanbul: Remzi Kitabevi.

\item Yeşilyaprak, B. (2016). \textit{21. yüzyılda eğitimde rehberlik hizmetleri} (26. baskı). Ankara: Nobel Yayınları.

\item YÖK. (2018). \textit{Psikolojik Danışmanlık ve Rehberlik Lisans Programı Ders İçerikleri}. Ankara: Yükseköğretim Kurulu.

\end{enumerate}

\newpage
